\chapter{Założenia i wymagania}
\label{chap:requirements}
\section{Punkt wyjścia: prototyp inżynierski}
\writingnote{Udokumentuj makietę, czujniki, napędy i poprzednią komunikację. Potwierdź model Raspberry Pi użyty wcześniej i odróżnij go od docelowej bramy badawczej.}

\section{Wymagania funkcjonalne}
\writingnote{Nadaj identyfikatory wymaganiom: obsługa niezależnych węzłów, zbieranie danych, status urządzeń, zapis historii, odczyt przez API, późniejsza detekcja i alerty. Powiąż wymagania z testami.}

\section{Wymagania niefunkcjonalne}
\writingnote{Określ mierzalne cele dotyczące opóźnienia, wykorzystania RAM, zachowania po awarii, odtwarzalności i pracy lokalnej. Nie wpisuj arbitralnych wartości jako uzyskanych wyników.}

\section{Granice zaufania i model zagrożeń}
\writingnote{Opisz zasoby, role i dostęp napastnika w laboratorium. Zdefiniuj, które zdarzenia system może obserwować i czego nie da się ustalić wyłącznie z telemetrii.}

\section{Zakres realizacji i ograniczenia}
Architektura ma umożliwiać pracę z jednym węzłem, typowym prototypem z 2--3 płytkami oraz dodatkowymi węzłami symulowanymi. Liczba fizycznych urządzeń nie stanowi stałej zapisanej w logice aplikacji.
\writingnote{Wyodrębnij etap bez elektroniki. TinyML, autoenkoder i rozbudowany dashboard pozostaw jako rozszerzenia opcjonalne.}
